% 主小行星带（综述）
% license CCBYSA3
% type Wiki

本文根据 CC-BY-SA 协议转载翻译自维基百科\href{https://en.wikipedia.org/wiki/Asteroid_belt}{相关文章}。

\begin{figure}[ht]
\centering
\includegraphics[width=6cm]{./figures/f9c18c4daeef0b22.png}
\caption{} \label{fig_belt_1}
\end{figure}
\begin{figure}[ht]
\centering
\includegraphics[width=6cm]{./figures/cf13395dc35061d2.png}
\caption{到目前为止，腰带内最大的物体是矮行星谷神星。小行星带的总质量明显小于冥王星大约是冥王星月亮的两倍卡隆。} \label{fig_belt_2}
\end{figure}
小行星带是太阳系中的一个圆环形区域，以太阳为中心，大致横跨行星木星和火星轨道之间的空间。该区域包含大量固态、不规则形状的天体，称为小行星或小行星状星体。这些已被确认的天体大小各异，但都远小于行星，且平均彼此相距约一百万千米（约六十万英里）。这个小行星带也被称为主小行星带或主带，以区别于太阳系中其他小行星族群$^{[1]}$。

小行星带是太阳系中最小且最内侧的环星盘。位于其他区域的小型太阳系天体类别，还包括近地天体、半人马天体、柯伊伯带天体、散射盘天体、塞德纳型天体以及奥尔特云天体。主带总质量中约 60\% 集中在四颗最大的小行星：谷神星（Ceres）、灶神星（Vesta）、智神星（Pallas）和健神星（Hygiea）中。小行星带的总质量估计约为月球质量的 3\%$^{[2]}$。

谷神星是小行星带中唯一足够大以被归类为矮行星的天体，其直径约为 950 km；而灶神星、智神星和健神星的平均直径均小于 600 km$^{[3][4][5][6]}$。其余在矿物学上已被分类的天体尺寸一直延续到仅数米大小$^{[7]}$。小行星带中的物质分布极为稀疏，以至于多艘无人航天器已多次穿越该区域而未发生任何碰撞事故$^{[8]}$。尽管如此，大型小行星之间的碰撞仍然会发生，并能形成小行星族群，其成员具有相似的轨道特征和成分。带内各个小行星依据其光谱进行分类，大多数可归入三大基本类型：富碳的 C 型、富硅酸盐的 S 型以及富金属的 M 型。

小行星带是由原始太阳星云中的一群微行星形成的$^{[9]}$，这些微行星是原行星的较小前身。然而，在火星与木星之间，来自木星的引力扰动阻碍了这些天体进一步聚积成一颗行星$^{[9][10]}$，并赋予它们过量的动能，从而在碰撞时将微行星以及大多数初生原行星击碎。其结果是，在太阳系历史最初的一亿年内，小行星带原始质量的 99.9\% 被消耗殆尽$^{[11]}$。其中一些碎片最终进入太阳系内部区域，导致与类地行星的陨石撞击。只要小行星绕太阳的公转周期与木星形成轨道共振，其轨道就会持续受到显著扰动；在这些轨道距离处会形成所谓的柯克伍德空隙，因为小行星被“扫”入其他轨道上$^{[12]}$。
\subsection{观测史}
\begin{figure}[ht]
\centering
\includegraphics[width=6cm]{./figures/afb56775c2129f99.png}
\caption{1596 年，约翰内斯·开普勒（Johannes Kepler）基于自己对于行星轨道比例关系的直觉，认为在火星与木星的轨道之间存在一颗尚未被发现的行星$^{[13]}$。} \label{fig_belt_3}
\end{figure}
1596 年，约翰内斯·开普勒（Johannes Kepler）在其著作《宇宙秘典》（Mysterium Cosmographicum）中写道：“在火星与木星之间，我放置一颗行星”，明确表达了他认为在那里应当存在一颗行星的预测$^{[14]}$。在分析第谷·布拉赫（Tycho Brahe）的观测数据时，开普勒认为火星与木星轨道之间的间隙过大，不符合他关于行星轨道应当分布位置的模型$^{[15]}$。

在 1766 年他匿名发表的对夏尔·博内（Charles Bonnet）《自然沉思录》（Contemplation de la Nature）的译本附注中$^{[16]}$，维滕贝格的天文学家约翰·丹尼尔·提丢斯（Johann Daniel Titius）$^{[17][18]}$ 指出行星布局中似乎存在某种规律，这一规律后来被称为提丢斯–波得定律（Titius–Bode Law）。如果从 0 开始构造一个数列，然后依次取 3、6、12、24、48 等数（每次加倍），再给每个数字加上 4 并除以 10，就可以得到与当时已知行星轨道半径（以天文单位计）惊人接近的数值——前提是必须承认在火星（对应数列的 12）与木星（48）轨道之间存在一颗“缺失的行星”（对应数列的 24）。在附注中，提丢斯写道：“难道造物主会把那片空间空出来吗？绝不会。”$^{[17]}$

当威廉·赫歇尔（William Herschel）于 1781 年发现天王星时，其轨道与这一定律高度吻合，使得一些天文学家断定，在火星和木星轨道之间必然存在一颗行星$^{[19]}$。
\begin{figure}[ht]
\centering
\includegraphics[width=6cm]{./figures/7ecae5c9889c484d.png}
\caption{朱塞佩·皮亚齐（Giuseppe Piazzi），即小行星带中最大天体谷神星（Ceres）的发现者：谷神星曾被视为一颗行星，但随后被重新分类为小行星，并自 2006 年起被归类为矮行星。} \label{fig_belt_4}
\end{figure}
1801 年 1 月 1 日，西西里巴勒莫大学天文学主席朱塞佩·皮亚齐（Giuseppe Piazzi）发现了一个微小的移动天体，其轨道半径与该规律所预测的数值完全一致。他将其命名为“谷神星”（Ceres），取自罗马的丰收女神、亦为西西里的守护神。皮亚齐起初认为它可能是一颗彗星，但其缺乏彗发的特征表明它更像是一颗行星$^{[20]}$。由此，上述规律成功预测了当时八颗行星（即水星、金星、地球、火星、谷神星、木星、土星和天王星）的轨道半长轴。

与谷神星的发现同时，应弗朗茨·克萨韦尔·冯·扎克（Franz Xaver von Zach）的邀请，24 名天文学家组成了一个非正式团体，被称为“天上警察”（celestial police），其明确目的在于寻找更多行星；他们重点在火星与木星之间搜索，因提丢斯–波得定律预言该区应存在一颗行星$^{[21][22]}$。

大约 15 个月后，“天上警察”成员海因里希·奥伯斯（Heinrich Olbers）在同一区域发现了第二个天体——智神星（Pallas）。但与其他已知行星不同，谷神星与智神星在望远镜最高放大倍率下仍然呈现为光点，而非解析为圆面。除了快速移动外，它们与恒星几乎无差别$^{[23]}$。

因此，1802 年，威廉·赫歇尔（William Herschel）建议将它们划入一个独立类别，并以希腊语 *asteroeides*（意为“类星的”）为来源，为其命名为“asteroids”（小行星）$^{[24][25]}$。在完成对谷神星与智神星的一系列观测后，他得出结论$^{[26]}$：

无论是“行星”还是“彗星”这一称谓，都不能恰当地应用于这两个天体……它们与小恒星极为相似，以至于几乎无法将其区分开来。基于它们这种类星的外观，我取其名，称之为“小行星”；但我保留更改其名称的权利，若日后发现一个能更好表达其本质的术语。

到 1807 年，进一步观测又在该区域发现了两颗新天体：婚神星（Juno）与灶神星（Vesta）$^{[23]}$。然而，拿破仑战争中利林塔尔（Lilienthal）天文台的焚毁使这一早期发现阶段被迫终结$^{[27][23]}$。

尽管赫歇尔已提出“小行星”一词，但在接下来的几十年里，人们仍普遍将这些天体称为行星$^{[16]}$，并按发现顺序为其编号：1 谷神星、2 智神星、3 婚神星、4 灶神星。然而 1845 年天文学家卡尔·路德维希·亨克（Karl Ludwig Hencke）发现了第五个天体（5 义神星 Astraea），紧接着更多天体以加速的速度被发现。将它们都算作行星变得愈发不切实际。最终，它们被从行星名单中移除（最早由亚历山大·冯·洪堡在 1850 年代初提出），而赫歇尔所造的“小行星”（asteroids）一词逐渐成为通用称谓$^{[16]}$。

1846 年海王星（Neptune）的发现使提丢斯–波得定律在科学界迅速失去信誉，因为其轨道与该定律所预测的位置毫不相近。直至今日，这一定律仍无科学解释，天文学界普遍认为它只是巧合$^{[28]}$。
\begin{figure}[ht]
\centering
\includegraphics[width=6cm]{./figures/922f458ab47f5d40.png}
\caption{951 号小行星加斯普拉（Gaspra）——首个被航天器成像的小行星，其图像来自伽利略号（Galileo）在 1991 年飞掠期间拍摄；图像的颜色经过了夸张处理。} \label{fig_belt_5}
\end{figure}
